\section{Introduction to Proofs [1.1-1.6]}

\subsection{Propositions}

\emph{A proposition is a statement, either true or false.}

\begin{itemize}
    \item Being T or F has limitations
    \item It excludes statements whose truth vary with circumstance
    \item It's not always easy to decide if a proposition is true or false
\end{itemize}

\vspace{1mm}

The symbol \(::=\) means 'equal by definition'.

\vspace{1mm}

You cannot check a claim about an infinite set, by checking a finite set of its elements.

\vspace{1mm}

Propositions about all numbers or items of some kind can be represented like this, eg \(\forall n \in \mathbb{N}. \; P(n) \; \text{is prime}\), where \(\forall\) is 'for all' and \(\mathbb{N}\) is the set of non-negative integers. 

\vspace{1mm}

For a computer scientist, some of the most important things are to prove the correctness of programs and systems, if they do what they are meant to.

\subsection{Predicates}

\emph{A predicate is a proposition whose truth depends on the value of one or more variables.}

\vspace{1mm}

You can use a function like notation to denote a predicate supplied with specific variable values.

\[P(n) ::= \text{'n is a perfect square'}\]

If \(P\) is a predicate, then \(P(n)\) is either true or false, depending on \(n\). However, if it was a function, then it would be an ordinary numerical quantity.

\subsection{The Axiomatic Method}

\emph{Axioms are propositions that are assumed to be true, a fundamental rule.}

You can start at these axioms, and establish the truth of further propositions by providing proofs. 

\vspace{1mm}

\emph{A proof is a sequence of logical deductions from axioms, and previously proved statements that concludes with the proposition in question.}

\emph{Theorems are important true propositions.}

\emph{A lemma is a preliminary proposition useful for proving later propositions. (Something that has already been proved, useful to prove new stuff.)}

\emph{A corollary is a proposition that follows in just a few logical steps from a theorem.}

\vspace{1mm}

Euclid's axiom and proof approach, called the axiomatic method, is the foundation of maths today.

A handful of axioms, the Zermelo-Fraenkel with Choice axioms (ZFC), with a few logical deduction rules, are sufficient to derive all of maths.

\subsection{Our Axioms}

ZFC is too primitive for practical purposes.

Be upfront about what you are proving, assuming and taking as result.

\subsubsection{Logical Deductions}

\emph{Logical deductions, or inference rules, are used to prove new propositions using previously proved ones.}

When statements above the line (\emph{antecedents}) are proved, we consider the statements below the line (\emph{the consequent}) to be proved.

\textbf{Modus Ponens - The law of detachment}

Modus Ponens is a fundamental inference rule that says a proof of \(P\) together with a proof of \(P \implies Q\) is a proof of \(Q\).

It is written:
\[\frac{P.\; P \implies Q}{Q}\]

or 

\[P \implies Q, \; P \vdash Q\]

If P implies Q, and P is true, therefore Q is true. The turnstile \(\vdash\) represents is a syntactic consequence, or is deducible from.

\textbf{Modus Tollens - The law of the Contrapositive}

\[\frac{Not(P) \implies Not(Q)}{Q \implies P}\]

or 

\[Q \implies P, \; \text{Hence} \;\neg P \implies \neg Q\] 

\subsubsection{Patterns of Proof}

There are several templates for proofs, which make it easier for you to start a proof. 

These are just outlines, so you always have something to say.

\subsection{Proving an Implication}

Implications are propositions in the form '\(\text{If} \; P \; \text{then} \; Q\)'. This is \( P \implies Q\).

There are standard methods to prove an implication.

\subsubsection{Method 1}

To prove \(P \implies Q\):

1) Write 'Assume \(P\)'

2) Show that \(Q\) logically follows.
\vspace{1mm}
Notes:

\begin{itemize}
    \item You might need to do scratchwork when starting a proof. This is fine, but not part of the proof. 
    \item Keep the proof clear and concise.
    \item Proofs typically start with '\emph{Proof}' and end with a delimiter eg QED. This clarifies where it starts and ends.
\end{itemize}

\subsubsection{Method 2 - Proving the Contrapositive}

An implication \(P \implies Q\) is logically equivalent to its contrapositive \( \neg Q \implies \neg P\).

\vspace{1mm}

The contrapositive is the inverse of the converse.

To prove \(\neg Q \implies \neg P\):

1) Write 'We prove the contrapositive', and state the contrapositive.

2) Proceed as in Method 1.

\subsection{Proving an If and Only If}

Two statements are logically equivalent - one holds if and only if the other holds.

\subsubsection{Method 1 - Prove each statement implies the other}

The statement \(P \iff Q\) is equivalent to \(P \implies Q\) and \(Q \implies P\).

1) Write 'We prove \(P \implies Q\) and vice versa.

2) Write, 'First we show \(P \implies Q\).' Do this with the methods from proving implications (assume p, show Q, or prove the contrapositive).

3) Write 'Now we show \(Q \implies P\).' Do this again with a method from proving implications.

\subsubsection{Method 2 - Construct a chain of Iffs}

To prove P is true iff Q is true:

1) Write 'We construct a chain of if-and-only-if implications.'

2) Prove P is equivalent to a second statement, equivalent to a 3rd etc until you reach Q.

\vspace{1mm}

This requires more ingenuity but the result can be a shorter proof.
